% Generated draft: requires mathematical and source review.
\documentclass[11pt]{article}
\usepackage[margin=25mm]{geometry}
\usepackage{fontspec}
\setmainfont{Latin Modern Roman}
\usepackage{amsmath,amssymb,mathtools,unicode-math}
\setmathfont{Latin Modern Math}
\usepackage{enumitem,needspace,xurl,hyperref}
\setlength{\parindent}{0pt}
\setlength{\parskip}{5pt}
\setlength{\emergencystretch}{3em}
\newcommand{\N}{\mathbb{N}}
\newcommand{\Z}{\mathbb{Z}}
\newcommand{\Q}{\mathbb{Q}}
\newcommand{\R}{\mathbb{R}}
\newcommand{\C}{\mathbb{C}}
\newcommand{\F}{\mathbb{F}}
\newcommand{\A}{\mathbb{A}}
\newcommand{\PP}{\mathbb{P}}
\newcommand{\E}{\mathbb{E}}
\newcommand{\eps}{\varepsilon}
\newcommand{\dd}{\,\mathrm{d}}
\begin{document}
\section*{UnsolvedMath ID: 800834}
\textbf{Source problem code:} catalogue entry
\subsection{Definitions and mathematical statement}
The supplied problem concerns bouncing droplets in certain vibrating liquids whose motion exhibits properties analogous to those of quantum mechanical particles. The liquid, forcing, droplet dynamics, and properties to be compared are not specified.

The named Madelung transform can be described locally as follows. For a complex scalar wave function admitting a sufficiently regular local phase representation away from its zeros, write
\[
\psi=\sqrt{\rho}\,\exp(iS/\hbar),\qquad \rho=|\psi|^2>0,
\]
where \(\hbar>0\) is the phase scale and \(S\) is real, defined modulo \(2\pi\hbar\). This is the usual amplitude--phase convention; it does not supply a wave equation or a droplet model. A phase need not exist globally, and this representation does not assign a phase at zeros of \(\psi\). An associated velocity field requires additional geometric and physical conventions not given in the record.

\textbf{Research problem (source formulation).} Explain the correspondence between bouncing droplets in certain vibrating liquids and quantum mechanical particles by combining averaging with the Madelung transform.

For this request to become a precise mathematical assertion, one must specify the droplet--liquid equations and their domain, forcing and boundary conditions; the averaging operation and its regime of validity; the target quantum equation and its state space; and the observables and sense in which the resulting descriptions correspond. None of these specifications is supplied. In particular, the source does not state an exact equivalence, a limiting theorem, an error estimate, or a quantifier over droplet systems or quantum states. The displayed transform defines the named operation locally, but does not resolve these missing specifications. The problem is therefore retained as an explicitly underdetermined research task.
\subsection{Short English statement}
Explain how averaging and the Madelung amplitude--phase transform could account for quantum-like behavior of bouncing droplets in certain vibrating liquids. The source leaves the physical model, averaging procedure, target quantum dynamics, and meaning of correspondence unspecified, so it does not determine a unique mathematical theorem.
\subsection{Sources}
\begin{itemize}
\item Local UnsolvedMath registry, permanent ID \texttt{800834}, source code \texttt{catalogue entry}, title: \emph{Geometric Hydrodynamics in Open Problems --- Problem 33}. The record attributes the problem to Khesin, Misiolek, and Modin, with proposed year 2022. Its status is supplied as \texttt{open}; this is a registry label, not an independently verified assessment.
\item Khesin, Misiolek, and Modin, \emph{Geometric Hydrodynamics in Open Problems} (2022), Problem 33. Supplied source URL: \url{https://arxiv.org/abs/2205.01143}. The source record reports access on 2026-10-16 and extraction by \texttt{arxiv\_latex}; these are supplied provenance statements, not a claim of independent consultation.
\item The problem text refers to ``[Couder, Bush]'' for hydrodynamic quantum analogues involving bouncing droplets. No titles, publication details, or links for these references are supplied, so their bibliographic identities cannot be fixed from this record.
\end{itemize}
\end{document}
